\documentclass{article}
\usepackage[a4paper, margin=1in]{geometry}
\usepackage{amsmath}
\usepackage{amssymb}
\usepackage{graphicx}
\usepackage{hyperref}
\usepackage{listings}

\title{Documentation for the Coupled Geometric PDE Solver}
\author{Gemini}
\date{\today}

\begin{document}

\maketitle

\section{Introduction}
This document details the implementation of a numerical solver for a system of coupled partial differential equations describing the evolution of a Riemannian metric $g$ and a one-form $\lambda$ on a triply periodic 3-dimensional manifold ($T^3$). The code, implemented in Python with NumPy, uses pseudospectral methods for spatial discretization and a 4th-order Runge-Kutta scheme for time integration.

The solver is contained in the file \texttt{coupled\_solver.py} and leverages several geometric helper functions to compute quantities such as the Lie derivative, Christoffel symbols, and the Laplace-de Rham operator.

\section{Mathematical Model}
The solver evolves a state $(g, \lambda)$ according to the following system of equations:

\begin{align}
    \partial_t g &= \mathcal{L}_v g \label{eq:g_evolution} \\
    \partial_t \lambda &= -dp - \nu \Delta_H \lambda \label{eq:lambda_evolution} \\
    \delta \lambda &= 0 \label{eq:constraint} \\
    v &= \lambda^\sharp \label{eq:closure}
\end{align}

Where the terms are defined as:
\begin{itemize}
    \item $g$: The Riemannian metric tensor field, $g_{ij}(x,t)$.
    \item $\lambda$: A time-dependent one-form, $\lambda_i(x,t)$.
    \item $v$: A time-dependent vector field, $v^i(x,t)$, related to $\lambda$ via the musical isomorphism (raising the index with the metric): $v^i = g^{ij}\lambda_j$.
    \item $\mathcal{L}_v g$: The Lie derivative of the metric $g$ with respect to the vector field $v$. In local coordinates, its components are given by:
    \begin{equation}
        (\mathcal{L}_v g)_{ij} = v^k \partial_k g_{ij} + g_{kj} \partial_i v^k + g_{ik} \partial_j v^k
    \end{equation}
    \item $\Delta_H$: The Laplace-de Rham operator (Hodge Laplacian). For a Ricci-flat metric ($R_{ij}=0$), this is equivalent to the negative of the connection Laplacian:
    \begin{equation}
        (\Delta_H \lambda)_i = -g^{jk}\nabla_j\nabla_k \lambda_i
    \end{equation}
    where $\nabla_k$ is the covariant derivative compatible with the metric $g$.
    \item $p$: A scalar field (pressure) that acts as a Lagrange multiplier to enforce the constraint.
    \item $dp$: The exterior derivative of $p$, which is the gradient of $p$ in this context.
    \item $\delta \lambda$: The codifferential of $\lambda$. The condition $\delta \lambda = 0$ is equivalent to the vector field $v$ being divergence-free: $\nabla_i v^i = 0$.
    \item $\nu$: A constant representing the viscosity.
\end{itemize}

\section{Numerical Implementation}
The system is solved in a triply periodic box of side length $L=2\pi$. The numerical method can be broken down into several key components.

\subsection{Time Integration}
A classic 4th-order Runge-Kutta (RK4) scheme is used to advance the state $(g, \lambda)$ in time. This involves evaluating the right-hand-side (RHS) of equations \eqref{eq:g_evolution} and \eqref{eq:lambda_evolution} at four different stages within each time step.

\subsection{Spatial Discretization}
All spatial derivatives are computed spectrally using the Fast Fourier Transform (FFT). This approach is highly accurate for smooth functions on a periodic domain. A function $\phi(x)$ is transformed to Fourier space $\hat{\phi}(k)$, multiplied by $ik$ (where $k$ is the wavenumber), and then transformed back to physical space to obtain its derivative.

\subsection{Pressure Projection}
To enforce the divergence-free constraint \eqref{eq:constraint} at each time step, a projection method is used. The constraint must hold for all time, which implies its time derivative must be zero: $\partial_t(\delta\lambda) = 0$.

A crucial point is that the time derivative $\partial_t$ and the codifferential $\delta$ do not commute, because $\delta$ depends on the time-varying metric $g$. Applying the condition $\partial_t(\delta\lambda)=0$ to the momentum equation \eqref{eq:lambda_evolution} yields a Poisson-Beltrami equation for the pressure $p$. The correct form for this equation is:
\begin{equation}
    \Delta_g p = -\delta(\nu \Delta_H \lambda) - \text{div}\left( (\mathcal{L}_v g)^\sharp \cdot v \right)
\end{equation}
where $\Delta_g p = \delta(dp)$ is the scalar Laplacian on $p$. The second term on the right-hand side is a source term arising from the time-dependence of the metric. It is written as the divergence of the vector field $W^i = g^{ij}(\mathcal{L}_v g)_{jk}v^k$. Because this source term is in divergence form, its integral over the periodic domain is guaranteed to be zero, thus satisfying the compatibility condition for the Poisson equation.

\subsection{Key Computational Functions}
\begin{description}
    \item[\texttt{compute\_lie\_derivative\_g}] Implements the formula for $\mathcal{L}_v g$.

    \item[\texttt{compute\_christoffel}] Computes the Christoffel symbols, $\Gamma^k_{ij}$, from the metric and its first derivatives.

    \item[\texttt{compute\_connection\_laplacian\_1form}] Computes $(\Delta_H \lambda)_i = -g^{jk}\nabla_j\nabla_k \lambda_i$, involving two nested covariant derivatives.

    \item[\texttt{BeltramiPoissonSolver.solve}] Solves the elliptic equation for the pressure $p$ given the full RHS, including the viscous and geometric source terms.
\end{description}

\section{Code Structure}
The solver is encapsulated within the \texttt{coupled\_solver.py} file.
\begin{itemize}
    \item \textbf{\texttt{CoupledSolver} class:} The main class that orchestrates the simulation. Its \texttt{solve} method contains the main RK4 time-stepping loop, and its \texttt{get\_rhs} method computes the time derivatives of $g$ and $\lambda$.
    \item \textbf{\texttt{BeltramiPoissonSolver} class:} A self-contained class for solving the scalar Poisson-Beltrami equation.
    \item \textbf{Helper Functions:} A suite of standalone functions for computing the necessary geometric quantities.
\end{itemize}

\end{document}
